% Source: https://artofproblemsolving.com/community/c3h1404927_yet_again_lithuanian_math_kangaroo_questions__student_2009
% Stable identity: aops-topic-1404927-q28
Among $1,\dots,99$ there are $33$ numbers of each residue $0,1,2$ modulo $3$. Two multiples of $3$ have a sum divisible by $3$, so each group contains at most one multiple of $3$, and therefore $n\ge33$.

$33$ groups are enough. A number $\equiv1$ and a number $\equiv2$ cannot share a group, but numbers $\equiv1$ may share a group with each other ($1+1\equiv2$), as may numbers $\equiv2$ ($2+2\equiv1$), and a multiple of $3$ may join either kind. So form $33$ groups, each containing one multiple of $3$. Put one number $\equiv1\pmod3$ into each of groups $1$ to $15$ and the remaining $18$ such numbers into group $16$; put one number $\equiv2\pmod3$ into each of groups $17$ to $32$ and the remaining $17$ into group $33$. Every group has at least $2$ numbers and no forbidden sum. The answer is $\boxed{\text{(C) } 33}$.
